\textbf{(CAP-COLUNI 2020 -- Q17)} O gráfico a seguir apresenta o crescimento das vendas de aparelhos celulares das marcas A e B no decorrer do tempo. O eixo vertical apresenta o número de milhões de aparelhos celulares vendidos e o eixo horizontal apresenta o intervalo de tempo considerado, de 2015 a 2018, no qual $x=1$ representa o ano de 2015; $x=2$, o ano de 2016 e, assim, sucessivamente.

De acordo com as informações apresentadas e mantendo o padrão de crescimento para as duas marcas de celular, é CORRETO afirmar que:

\includegraphics{./figs/2020-Q17.png}

\begin{enumerate}
	\item a função $y=10x+26$ descreve o número de aparelhos vendidos da marca B ao longo dos anos.
	\item em 2020, serão vendidos 104 milhões de aparelhos da marca B.
	\item o número de aparelhos vendidos pelas duas marcas será o mesmo apenas no ano de 2021.
	\item no ano de 2017, a diferença entre o número de aparelhos vendidos das marcas A e B foi de 35 milhões.
\end{enumerate}
